\section{2025 Forecast}

Owing to the absence of credible environmental factors data for 2025, we have estimated the environmental and human factors from our dataset. Then we have applied the Gaussian Ridge Regression for crop yield prediction model and Bayesian Optimization for Crop Yield Maximization Model to get the suggested human factors for optimal Crop Yield.  

\subsection{Predicted Values for 2025 using Time Series and Forecasting Crop yield}

To predict crop yield for the year 2025, we require future estimates of the relevant predictor variables. Each of these variables is a time series and exhibits a trend over the years. To obtain the 2025 forecasts, we applied the following modeling approaches:

\begin{itemize}
    \item \textbf{Quadratic Trend Fitting:} Applied to continuous predictor variables — Rainfall, Average Temperature, Fertilizer Use, and Technology Index — to capture any nonlinear time-based patterns.
    \item \textbf{Poisson Regression:} Applied to the count-based variable — Pest Incidents — as it effectively models count data and event rates over time.
\end{itemize}

\begin{center}
\textbf{\large 2025 FORECASTS (Polynomial Trends)}

\vspace{0.5cm}

\begin{tabular}{>{\raggedright}p{5cm}c}
\toprule
\textbf{Variable} & \textbf{Forecasted Value} \\ 
\midrule
Rainfall & 1236.78 mm \\
Average Temperature & 25.82 °C \\
Fertilizer Use & 261.0338 kg/ha \\
Technology Index & 74.07 \\
Pest Incidents & 3 events \\
\bottomrule
\end{tabular}
\end{center}


Using our optimized Random Forest model, we have predicted the crop yield for the specified agricultural conditions in 2025 is \textbf{7.0971 tons/ha}


The prediction of 7.0971 tons/ha suggests:

\begin{itemize}
\item \textbf{Above-Average Performance}: This yield is 6.61 $\%$ higher than the dataset mean of 6.66 tons/ha
\item \textbf{Key Contributing Factors}:
  \begin{itemize}
  \item High Technology Index (74.07) likely contributed significantly
  \item Optimal fertilizer application (261.03 kg/ha) supported growth
  \item Moderate rainfall (1236.78 mm) provided adequate moisture
  \end{itemize}
\item \textbf{Potential Limitations}:
  \begin{itemize}
  \item Pest incidents (3 events) may be slightly constraining yields
  \item Temperature (25.82°C) is near [optimal/suboptimal] range
  \end{itemize}
\end{itemize}

\subsection{Predicting Crop Yield through the Optimization Problem}

The following results are obtained by using the Bayesian Optimization process though the Random Forest Model.

\subsubsection{Input Conditions}
The optimization was performed with the following environmental conditions:

\begin{table}[H]
\centering
\caption{Current Environmental Conditions}
\label{tab:conditions}
\begin{tabular}{@{}lc@{}}
\toprule
\textbf{Variable} & \textbf{Value (Unit)} \\
\midrule
Rainfall & 1236.78 mm \\
Average Temperature & 25.82 \si{\degreeCelsius} \\
Pest Incidents & 3 count \\
\bottomrule
\end{tabular}
\end{table}

\subsubsection{Optimization Results}
The Bayesian optimization process produced the following optimal resource allocation within the \SI{1500}{units} budget constraint:

\begin{table}[H]
\centering
\caption{Optimal Resource Allocation}
\label{tab:results}
\begin{tabular}{@{}lc@{}}
\toprule
\textbf{Parameter} & \textbf{Value (Unit)} \\
\midrule
Fertilizer Application & 312.34 \si{kg/ha} \\
Technology Investment & 61.48 units \\
Predicted Crop Yield & 7.78 \si{ton/ha} \\
Total Expenditure & 932.08 units \\
Budget Utilization & 93.2\% \\
\bottomrule
\end{tabular}
\end{table}

\subsubsection{Yield Improvement Analysis}
The optimization achieved a predicted yield of \textbf{7.78 ton/ha}, representing a \textbf{9.7\% increase} over the original forecast of 7.09 ton/ha without optimization. This improvement is statistically and practically significant because:

\begin{itemize}
\item The 0.69 ton/ha increase would translate to substantial economic gains at scale (approximately 690 kg additional yield per hectare)
\item The improvement was achieved while maintaining 93.2\% budget utilization, demonstrating efficient resource allocation
\item The yield increase comes without additional environmental risk, as fertilizer application remains within reasonable agronomic limits
\end{itemize}

The results demonstrate that Bayesian optimization can effectively balance multiple constraints (budget, environmental conditions) while identifying non-obvious optimal combinations of inputs that outperform conventional allocation strategies. The near-full budget utilization (93.2\%) suggests the algorithm successfully identified the point of diminishing returns for additional investments.

\subsubsection{Optimization Process Visualization}

\begin{figure}[h]
\centering
\begin{subfigure}{0.48\textwidth}
\includegraphics[width=\linewidth]{yield optimization progress.jpeg}
\caption{Yield Optimization Progress}
\label{fig:yield}
\end{subfigure}
\hfill
\caption{Bayesian Optimization Convergence}
\label{fig:optimization}
\end{figure}

\subsubsection*{Figure Explanations}

\textbf{Figure \ref{fig:yield} (Yield Optimization Progress):}
\begin{itemize}
\item Shows the evolution of predicted crop yield across optimization iterations
\item The algorithm quickly finds solutions with yields around 7.78 ton/ha
\item The somewhat stable line after iteration 40 indicates convergence to the optimal solution
\item Negative values early on represent penalty scores for budget violations
\item The later fluctuations indicates the algorithm trying new combinations. Bad results indicate the respective combination is discarded in the next iterations.
\end{itemize}



\subsubsection{Summary}
The optimization achieved a predicted yield of $7.78~\text{ton/ha}$ while utilizing 93.2\% of the available budget. Key findings include:
\begin{itemize}
\item The model recommended applying $312.34~\text{kg/ha}$ of fertilizer
\item Suggested technology investment of 61.48 units
\item The solution leaves 67.92 units of the budget unspent
\item Convergence plots demonstrate the optimization process effectively explored the parameter space
\end{itemize}

These recommendations are based on current conditions of approximately \SI{1236.78}{\mm} rainfall and average temperature of \SI{25.82}{\degreeCelsius} with 3 pest incidents.


